% HD-F1-T33. GPT-6 Astra Pro, /root/transport; 22 November 2026 UTC.
\section[HD-F1-T33: finite Fourier exclusion of linear level one]{HD-F1-T33: the actual finite Fourier model excludes every remaining linear level-one source over fields of size two and three}
\label{sec:f1-linear-level-one-fourier-exclusion}

No remaining T2-OTHER source has projective socle
\(PSL_N(q)\), with \(q\in\{2,3\}\), \(N\ge6\),
and level-one constituents on the actual perfect central
layer. The bound applies to every central and outer lift
and to all actual restriction multiplicities. Together
with T32, it gives the effective necessary bound
\[
 \boxed{\quad N<10^{62}\quad}
 \qquad\text{for every remaining }PSL_N(q)\text{ source.}
 \tag{F1T33.1}
\]
For the unitary family the retained T32 alternative
remains \(N<10^{62}\), or level one over \(q=2,3\).
The present proof establishes its conclusion directly
on the original coefficient space.

\paragraph{1. Exact scalar sectors and irreducibility.}
Let \(A=\mathbf F_q^N\), \(\mathcal H=\mathbf C[A]\),
with its orthonormal delta-function basis, and let
\[
 (P_a f)(x)=f(a^{-1}x)\qquad(a\in GL_N(q)).
 \tag{F1T33.2}
\]
The trace of \(P_h\) is \(q^{\dim\ker(h-1)}\), so
\(\mathcal H\) affords exactly the defining reducible
Weil character used to define level in T32.
For a character \(\eta\) of \(\mathbf F_q^\times\),
put
\[
 p_\eta=\frac1{q-1}\sum_{s\in\mathbf F_q^\times}
                  \eta(s)^{-1}P_{sI},
 \qquad E_\eta=p_\eta\mathcal H.
 \tag{F1T33.3}
\]
Write \(U=\langle\delta_0,\mathbf1_A\rangle\), a
two-dimensional invariant space. Define
\[
 W_1=E_1\cap U^\perp,\qquad
 W_\eta=E_\eta\quad(\eta\ne1).
 \tag{F1T33.4}
\]
The scalar action is free on each nonzero line, hence
\[
 d_\eta:=\dim W_\eta
 =\frac{q^N-1}{q-1}-\mathbf1_{\eta=1},
 \qquad
 \mathcal H=U\oplus\bigoplus_\eta W_\eta.
 \tag{F1T33.5}
\]

Here every \(W_\eta\) is irreducible for
\(H=SL_N(q)\), and different sectors are inequivalent.
To prove this, consider the transitive action of \(H\)
on \(A\setminus\{0\}\). The stabilizer of \(e_1\)
has exactly \(q\) orbits: each of the \(q-1\)
nonzero scalar multiples of \(e_1\) is fixed, and all
vectors outside \(\mathbf F_qe_1\) form one orbit.
For the last assertion, extend \((e_1,y)\) and
\((e_1,y')\) to bases. A map between the bases can
have its determinant adjusted to one by rescaling a
third basis vector; this uses \(N\ge3\).

An equivariant endomorphism of the permutation module
\(\mathbf C[A\setminus\{0\}]\) has matrix entries
constant on the orbits of ordered pairs. Therefore
its endomorphism algebra has dimension exactly \(q\).
Its scalar decomposition contains the trivial line,
the nonzero augmentation submodule \(W_1\), and the
\(q-2\) nontrivial scalar sectors \(W_\eta\).
These are \(q\) nonzero invariant direct summands.
Semisimplicity shows that any further reducibility,
or any isomorphism between distinct summands, would
increase the dimension of the endomorphism algebra
above \(q\): an irreducible constituent of multiplicity
\(r\) contributes \(r^2\). Thus the stated
irreducibility and inequivalence hold.
It follows from the least-power definition of level
that the \(W_\eta\) are precisely all level-one
irreducible \(H\)-modules; the two summands in \(U\)
have level zero.

\paragraph{2. The literal finite Fourier intertwiner.}
Since \(q\) is prime, fix
\(\psi(t)=\exp(2\pi i t/q)\). Define the actual unitary
operator
\[
 (\mathcal F f)(x)=q^{-N/2}
       \sum_{y\in A}\psi(x^t y)f(y).
 \tag{F1T33.6}
\]
Additive-character orthogonality gives
\(\mathcal F^2=P_{-I}\) and
\(\mathcal F^*\mathcal F=1\). A change of variable
\(y=az\) in this finite sum gives
\[
       \mathcal F P_a\mathcal F^{-1}=P_{a^{-t}}.
 \tag{F1T33.7}
\]
Also \(\mathcal F P_{sI}=P_{s^{-1}I}\mathcal F\).
Every character of \(\mathbf F_q^\times\) is its own
inverse when \(q=2,3\), so every \(E_\eta\) is
\(\mathcal F\)-stable. Moreover
\[
 \mathcal F\delta_0=q^{-N/2}\mathbf1_A,
 \qquad
 \mathcal F\mathbf1_A=q^{N/2}\delta_0.
 \tag{F1T33.8}
\]
Thus \(U\), its orthogonal complement, and every
\(W_\eta\) are stable. Equations (F1T33.7)--(F1T33.8)
are the exact intertwiners for the graph automorphism
and the removal of the invariant two-space.

For these prime fields, Malle--Testerman
Theorem~24.24 has no nontrivial field-automorphism
factor. Every automorphism of \(PSL_N(q)\) is therefore
represented on the universal cover by
\[
 h\longmapsto aha^{-1}
 \quad\text{or}\quad
 h\longmapsto ah^{-t}a^{-1},
              \qquad a\in GL_N(q).
 \tag{F1T33.9}
\]
The universal-cover uniqueness square of T32 identifies
these standard automorphisms with the actual lift of
each original \(g\in G\). The implementing operator
on each \(W_\eta\) is exactly
\(P_a|_{W_\eta}\) or
\((P_a\mathcal F)|_{W_\eta}\).
In particular every actual automorphism stabilizes
each chosen irreducible constituent, not merely the
sum of the scalar sectors.

\paragraph{3. Direct trace bounds on the same modules.}
First let \(a\in GL_N(q)\) be nonscalar.
For \(q=2\), writing \(d=\dim\ker(a-1)\le N-1\)
and removing \(U\) gives
\[
 \operatorname{Tr}_{W_1}(P_a)=2^d-2,
 \qquad
 |\operatorname{Tr}_{W_1}(P_a)|
 \le2^{N-1}-2<\frac{d_1}{2}.
 \tag{F1T33.10}
\]
Here \(N\ge3\) ensures that the positive upper bound
also dominates the possible absolute value one at
\(d=0\).

For \(q=3\), put
\(d_+=\dim\ker(a-1)\) and
\(d_-=\dim\ker(a+1)\). Both are at most \(N-1\),
and their eigenspaces have direct sum, so
\(d_++d_-\le N\). From the actual projectors
(F1T33.3),
\[
 \begin{aligned}
 \operatorname{Tr}_{W_1}(P_a)
     &=\frac{3^{d_+}+3^{d_-}}2-2,\\
 \operatorname{Tr}_{W_{-}}(P_a)
     &=\frac{3^{d_+}-3^{d_-}}2,
 \end{aligned}
 \tag{F1T33.11}
\]
where \(W_-\) is the sector of the nontrivial scalar
character. The inequalities on \(d_\pm\) imply
\(3^{d_+}+3^{d_-}\le3^{N-1}+3\).
Indeed, order the two exponents. When the smaller one
is positive, replacing \((u,v)\) by \((u+1,v-1)\)
increases the sum, and one may continue until the
smaller exponent is one, or the larger is \(N-1\).
The case with smaller exponent zero is even smaller.
The lower value of the first expression in
(F1T33.11) is at least \(-1\). Consequently
\[
 |\operatorname{Tr}_{W_\eta}(P_a)|
 \le\frac{3^{N-1}-1}{2}
 \le\frac{d_\eta}{3}
       \qquad(q=3).
 \tag{F1T33.12}
\]

For the graph coset take any \(a\in GL_N(q)\).
The diagonal entries of \(P_a\mathcal F\) in the
delta basis all have absolute value \(q^{-N/2}\).
Therefore
\[
 |\operatorname{Tr}_{\mathcal H}(P_a\mathcal F)|
       \le q^{N/2}.
 \tag{F1T33.13}
\]
Inserting the projector (F1T33.3), each term in its
average is another operator \(P_{a'}\mathcal F\).
Hence the same bound holds on \(E_\eta\).
The restriction to the unitary invariant space \(U\)
has trace of absolute value at most two, so
\[
 |\operatorname{Tr}_{W_\eta}(P_a\mathcal F)|
       \le q^{N/2}+2\le\frac{d_\eta}{2}
                  \qquad(N\ge6).
 \tag{F1T33.14}
\]
Here is an integer check of the last uniform inequality.
Use \(d_\eta\ge(q^N-q)/(q-1)\) and
\(q^{N/2}\le q^{N-3}\). The difference between the
required right and left sides after clearing
\(2(q-1)\) is at least
\[
       (q^3-2q+2)q^{N-3}-(5q-4).
 \tag{F1T33.15}
\]
It is at least \(42\) for \(q=2\), and at least
\(610\) for \(q=3\), since \(N\ge6\).
Thus every nonscalar inner representative and every
graph-coset representative has trace of absolute
value at most one half of the dimension on every
relevant irreducible module.

\paragraph{4. Exact factorization of every original matrix, including phases.}
Keep the original \(G,V,Z,L\) and the actual
surjection \(\beta:H\to L\) of T32. Universal-cover
exceptions for these linear families have \(N\le4\),
so this map is valid throughout \(N\ge6\).
The irreducible constituents of \(V_{\mathbf C}|_H\)
are level one by the present source branch. Paragraph~2
shows that every original conjugation fixes each such
constituent. The original irreducibility and the
isotypic-orbit argument of T32 therefore make the
restriction isotypic for one \(W=W_\eta\).
The evaluation map is an actual isomorphism
\[
 \operatorname{ev}:W\otimes M\xrightarrow{\ \sim\ }
       V_{\mathbf C},\qquad
 M=\operatorname{Hom}_H(W,V_{\mathbf C}),\qquad
 w\otimes f\longmapsto f(w).
 \tag{F1T33.16}
\]

For every \(g\in G\), choose the unitary implementing
operator \(A_g=P_a\) or \(P_a\mathcal F\) supplied
by (F1T33.9), restricted to this \(W\). Define
\[
 T_g(f)=\rho(g)\circ f\circ A_g^{-1}.
 \tag{F1T33.17}
\]
The commuting square of T32 makes \(T_g(f)\)
\(H\)-equivariant. Evaluation gives the precise equality
\[
       \rho(g)\operatorname{ev}
       =\operatorname{ev}(A_g\otimes T_g).
 \tag{F1T33.18}
\]
This retains every original scalar phase in \(T_g\).
No multiplicity-one assertion or independent
product-one normalization of these factors is used.

To see the trace estimate on the multiplicity factor,
average a positive Hermitian form on \(V_{\mathbf C}\)
over the finite original group. With the standard
form on \(W\), Schur's lemma makes
\(f^*f'\) a scalar operator, and these scalars define
a positive Hermitian form on \(M\). Evaluation is
then an isometry for the tensor product form.
Since \(\rho(g)\) and \(A_g\) are unitary,
(F1T33.18) makes \(T_g\) unitary as well. Hence
\(|\operatorname{Tr}(T_g)|\le\dim M\).
For every nonscalar original \(g\), its projective
automorphism is nontrivial. An inner representative
in (F1T33.9) is therefore nonscalar, and the graph
case is covered for every \(a\). Equations
(F1T33.10)--(F1T33.18) give
\[
 |\operatorname{Tr}_V(g)|
 =|\operatorname{Tr}_W(A_g)\operatorname{Tr}_M(T_g)|
 \le\frac12\dim W\dim M=\frac m2.
 \tag{F1T33.19}
\]
Every eigenvalue of \(\rho(g)\) has absolute value
one, so
\[
              \operatorname{rank}(g-1)\ge\frac m4.
 \tag{F1T33.20}
\]
A nonidentity scalar original inertia has moving
rank \(m\), and satisfies the same lower bound.
Finally the original source budget gives
\[
 2m+6=\sum_{\nu=1}^b\operatorname{rank}(g_\nu-1)
          \ge\frac{bm}{4}\ge3m,
 \qquad m\le6.
 \tag{F1T33.21}
\]
This contradicts \(m\ge27\) in the remaining gate,
and proves the exclusion and its consequence
(F1T33.1).

\paragraph{5. Actual source and realization maps.}
The finite Fourier operator, scalar projectors,
evaluation and multiplicity operators are the explicit
coefficient-space maps (F1T33.3), (F1T33.6),
(F1T33.16)--(F1T33.18). They estimate the same
original inertia matrices; no auxiliary cover is used.
The map on the original curve cover is its identity,
as are the maps on the coefficient middle extension,
branch stalks, parabolic cup form and prescribed
sixfold primitive corner. The rational projector,
its integer-cleared Jacobian image and the restricted
theta polarization therefore retain their identity
maps and multiplier one. These induce the identity
maps on the Poincare biextension, Hodge filtration,
period differential, and Grothendieck/Adams
realizations of F01--F04. The result excludes an
actual source family and adds no detector component.
